\chapter{Reference: errors and recovery}
\label{ch:refman-errors}

\begin{aside}
Things will go wrong. Files won't open. Network calls will
fail. Subprocesses will exit nonzero. Handle every failure
in \code{update}; show the user something humane.
\end{aside}

\section{Working example}

\begin{lstlisting}[language=C++]
#include <maya/maya.hpp>
#include <fstream>
#include <sstream>

using namespace maya;
using namespace maya::dsl;

struct Errors {
    struct Model {
        std::string content;
        std::string error;
    };
    struct LoadOk { std::string body; };
    struct LoadErr { std::string msg; };
    struct Load {};
    struct Quit {};
    using Msg = std::variant<Load, LoadOk, LoadErr, Quit>;

    static Model init() { return {}; }

    static auto update(Model m, Msg msg)
        -> std::pair<Model, Cmd<Msg>>
    {
        return std::visit(overload{
            [&](Load) {
                auto task = []() -> std::variant<std::string, std::string> {
                    std::ifstream f("/etc/hostname");
                    if (!f) return std::variant<std::string, std::string>{
                        std::in_place_index<1>,
                        "could not open /etc/hostname"
                    };
                    std::stringstream ss;
                    ss << f.rdbuf();
                    return std::variant<std::string, std::string>{
                        std::in_place_index<0>, ss.str()
                    };
                };
                return std::pair{m, Cmd<Msg>::task(task,
                    [](auto v) -> Msg {
                        if (v.index() == 0) return LoadOk{std::get<0>(v)};
                        return LoadErr{std::get<1>(v)};
                    })};
            },
            [&](LoadOk ok) {
                m.content = ok.body;
                m.error.clear();
                return std::pair{m, Cmd<Msg>{}};
            },
            [&](LoadErr e) {
                m.error = e.msg;
                return std::pair{m, Cmd<Msg>{}};
            },
            [](Quit) {
                return std::pair{Model{}, Cmd<Msg>::quit()};
            },
        }, msg);
    }

    static Element view(const Model& m) {
        return v(
            t<"file viewer"> | Bold,
            blank_,
            when(!m.error.empty(),
                v(
                    t<"error"> | Bold | Fg<255, 80, 80>,
                    text(m.error) | Fg<255, 180, 180>
                ),
                text(m.content)),
            blank_,
            t<"l load, q quit"> | Dim
        ) | pad<1> | border_<Round>;
    }

    static auto subscribe(const Model&) -> Sub<Msg> {
        return key_map<Msg>({
            {'q', Quit{}},
            {'l', Load{}},
        });
    }
};

int main() { run<Errors>({.title = "errors"}); }
\end{lstlisting}

\section{Notes}

\begin{itemize}[leftmargin=*]
  \item Errors are data; carry them in the Model just like
    any other field.
  \item Don't throw across the task boundary; encode failures
    as variants and dispatch.
  \item Provide a retry path; a stuck error message is
    almost as bad as a crash.
\end{itemize}

\section{Recap}

\begin{recap}
\begin{itemize}[leftmargin=*]
  \item Errors are messages.
  \item Variant of OK / err results from tasks.
  \item Render the error in the same view; offer retry.
\end{itemize}
\end{recap}

\section*{Workshop}

\begin{enumerate}[leftmargin=*]
  \item Add a retry key when an error is showing.
  \item Auto-clear errors after five seconds.
  \item Log errors to a file alongside the UI.
\end{enumerate}
